% ============================================================
%  Chapter 5 — Results and Refinements
%  Master's Thesis: Evaluating and Refining Hypo-Stage
%  José Gonçalves Lima Neto — IME-USP, 2026
% ============================================================
\chapter{Results and Refinements}
\label{chap:results}
\label{ch:results}

This chapter presents the empirical findings of this study and the refinements applied to
Hypo-Stage that emerged from them. The results are organized around the four
research questions defined in Chapter~\ref{chap:introduction}: 
\begin{itemize}
  \item RQ1 addresses how teams used the tool and the technique in practice;
  \item RQ2 characterizes the architectural uncertainties they registered;
  \item RQ3 reports the benefits and challenges perceived by participants;
  \item RQ4 documents each refinement, tracing it to a concrete observation and to the question it addresses.
\end{itemize}

The section~\ref{sec:emerging-patterns} goes beyond a direct answer to those four
questions by drawing two families of emerging findings: a set of
\textit{Adoption-Condition Candidates} --- empirically grounded candidate
conditions (not validated patterns) observed to shape whether hypothesis
engineering becomes a sustainable practice --- and a set of novel
\textit{Structural Hypothesis Patterns} for hypothesis engineers. The Structural
Hypothesis Patterns constitute a distinct contribution of this work, addressing a
gap left open by prior research in the ArchHypo line, as discussed in
Section~\ref{sec:group-b-patterns}. The Adoption-Condition Candidates describe
problem and context only; the derivation of their solutions is deferred to future
work.

\section{Research Methodology Traceability}
\label{sec:methodology-traceability}

The evidence presented in this chapter was collected and analyzed strictly in
accordance with the research design described in Chapter~\ref{chap:research-design}.
This section makes that correspondence explicit so that compliance with the
research methodology can be assessed at a glance; it does not restate the
methodology itself, which is defined in Chapter~\ref{chap:research-design}.

The evidence draws on the three data sources specified in
Section~\ref{sec:data-collection}: participant observation, artifact analysis,
and open-ended active-listening sessions. The quantitative summaries in
Sections~\ref{sec:rq2-results} and~\ref{subsec:refinement-summary} are derived
exclusively from the Hypo-Stage artifact records. As explained in
Section~\ref{subsec:active-listening}, the active-listening sessions were not
recorded; paraphrased, role-attributed observations drawn from the researcher's
reconstructed notes are reported only when they add evidence beyond the
recorded artifacts.

These data were analyzed with the three-stage thematic coding procedure defined
in Section~\ref{sec:data-analysis}. Table~\ref{tab:code-traceability} provides
the local traceability that this analysis adds: it lists representative open
codes and traces each one to its axial category, its research question (or
emerging-pattern group), and the chapter section where the corresponding finding
is developed.

\begin{table}[htbp]
\caption{From open codes to chapter sections: traceability of the thematic coding process.}
\label{tab:code-traceability}
\centering\small
\begin{tabular}{p{0.26\linewidth} p{0.26\linewidth} l p{0.17\linewidth}}
\toprule
\textbf{Open code} & \textbf{Axial category} & \textbf{RQ / Group} & \textbf{Discussed in Section} \\
\midrule
\textit{kafka-cluster-burst-registration}
  & Patterns of tool use; Structural archetypes
  & RQ1; RQ2
  & \ref{sec:results-rq1}, \ref{sec:group-b-patterns}~(B1) \\
\midrule
\textit{senior-author-only}
  & Organizational adoption conditions
  & Adoption candidate
  & \ref{sec:group-a-patterns}~(A1) \\
\midrule
\textit{registration-barrier-time-pressure}
  & Perceived challenges (registration time cost)
  & RQ3
  & \ref{subsec:challenges} \\
\midrule
\textit{tool-positive-catalog-integration}
  & Perceived benefits (catalog / Backstage integration)
  & RQ3
  & \ref{subsec:benefits} \\
\midrule
\textit{falsifiability-confusion}
  & Perceived challenges (negative-formulation learning curve); Candidate refinements
  & RQ3; RQ4
  & \ref{subsec:challenges}, \ref{sec:refinements} \\
\bottomrule
\end{tabular}
\end{table}

The artifact data from Hypo-Stage were summarized descriptively (hypothesis
counts, quality-attribute frequency, technical plan distribution) to complement
and illustrate the qualitative findings. Following methodological guidance for flexible qualitative designs \citep{robson2016}, the
inclusion of these descriptive statistics is consistent with a flexible
qualitative design and does not constitute a multistrategy approach.

% ─────────────────────────────────────────────────────────────
\section{RQ1 — How Teams Used Hypo-Stage in Practice}
\label{sec:rq1-results}
\label{sec:results-rq1}
% ─────────────────────────────────────────────────────────────

Unless noted otherwise, the empirical findings for RQ1--RQ3 in
Sections~\ref{sec:rq1-results}--\ref{sec:rq3-results} describe the full
observation window cumulatively; the refinements documented under RQ4 were
deployed iteratively during that same window, in the
observe--refine--deploy cycle described in Section~\ref{sec:procedure}, so
the results below are not restricted to the pre-refinement tool state.

\subsection{Hypothesis Registration Dynamics}
\label{subsec:registration-dynamics}

Over the observation period (Hypo-Stage registrations span 11~February~2026--12~April~2026, approximately nine calendar weeks or five two-week sprints), both teams together registered thirteen
architectural hypotheses in Hypo-Stage. Six hypotheses originated from the
Product~Team and seven from the Platform~Team. This distribution is noteworthy
because the Platform~Team is significantly smaller (four engineers) yet produced
more registered hypotheses than the twelve-engineer Product~Team. This asymmetry
is partly explained by the registration of a dense cluster of five Platform
hypotheses on a single day (12 April 2026), all addressing complementary aspects
of the shared Kafka streaming infrastructure — a burst pattern discussed further
in Section~\ref{sec:group-b-patterns}.

Across both teams, twelve of the thirteen hypotheses were authored by senior
technical leaders — senior or staff engineers who held explicit architectural
responsibilities for their respective domains. Table~\ref{tab:case-comparison}
reports the underlying team composition: five senior and seven mid-level
engineers on the Product~Team, and two senior and two mid-level engineers on
the Platform~Team. One hypothesis was registered by a mid-level engineer on the
Product~Team. Both technical leaders in each team
independently reported that a much larger number of architectural uncertainties
existed in their daily work than the number they managed to register. The
principal constraint they named was time: formulating a well-structured hypothesis
statement requires focused, uninterrupted attention that sprint delivery
pressure rarely afforded. This observation resonates with findings from prior
ArchHypo evaluations, where the learning curve and process-integration demands
were described as significant barriers to adoption
\citep{silva2024tse, silva2024ieeesw}.

\subsection{Role-Based Engagement}
\label{subsec:role-engagement}

Participant observation and the active listening sessions revealed a clear
stratification in how different roles engaged with Hypo-Stage during the study
period. Senior engineers treated the tool as a decision-support instrument:
they used it to externalize concerns that had until then existed only as tacit
knowledge, and they referenced the hypothesis backlog when discussing
architectural trade-offs in planning sessions. Mid-level and junior engineers
were aware of the tool's existence and attended the training session, but
neither team saw them register hypotheses independently beyond the initial
introduction. This does not indicate disengagement; rather, both technical
leaders explained that these engineers typically operate at a scope where
architectural uncertainties are resolved quickly — within a sprint or less —
leaving little perceived need for formal registration.

\subsection{Tool Adoption Patterns}
\label{subsec:adoption-patterns}

Usage of Hypo-Stage was non-uniform across the observation window. Registration
activity was concentrated in two periods: a first wave during and immediately
after the training session, and a second, more intensive wave near the end of
the period when the Platform~Team used the tool for a structured architectural
discussion on Kafka failure-handling strategies. Between these waves, both
technical leaders reported using the hypothesis detail view to review existing
records and update technical plans, but the tool was not used in team-visible
ceremonies such as sprint retrospectives or planning meetings. When asked,
participants noted that embedding the tool into these ceremonies would require
an explicit cultural agreement within the team — something that had not yet
occurred within the observation window.

The integration of Hypo-Stage into the Backstage Internal Developer Portal was
cited by both teams as a meaningful enabler of adoption. Because Backstage is
already the daily entry point for service catalog, documentation, and CI/CD
status, engineers encountered Hypo-Stage without navigating to a separate tool.
The EntityHypothesesTab, shown in Figure~\ref{fig:entity-hypotheses-tab}, which surfaces only the hypotheses linked to the
component an engineer is already viewing, was specifically mentioned as reducing
the cognitive cost of context-switching into hypothesis management.

\begin{figure}[ht]
  \centering
  \includegraphics[width=0.85\textwidth]{entity-hypotheses-tab}
  \caption{EntityHypothesesTab: Hypo-Stage surfaces only the hypotheses linked
    to the Backstage catalog component the engineer is currently viewing,
    eliminating the need to navigate to a separate tool or filter by component.}
  \label{fig:entity-hypotheses-tab}
\end{figure}

% ─────────────────────────────────────────────────────────────
\section{RQ2 — Architectural Uncertainties Registered}
\label{sec:rq2-results}
\label{sec:results-rq2}
% ─────────────────────────────────────────────────────────────

\subsection{Data Overview}
\label{subsec:data-overview}

Table~\ref{tab:hypotheses-overview} presents the thirteen hypotheses registered
during the study, anonymized in accordance with the ethical agreement described
in Section~\ref{sec:ethical-considerations}. Full exported records (statements,
ratings, entity links, and author notes) appear in Appendix~\ref{app:hypotheses-data}
(Section~\ref{sec:app-hypothesis-records}, \textit{Hypothesis Records}). Each hypothesis is identified by a
short label, its source type (Requirements or Solution), its lifecycle status
at the end of the observation period, and its uncertainty and impact ratings.

\begin{table}[htbp]
\caption{Overview of the thirteen registered architectural hypotheses.
  Req.\ = Requirements-sourced; Sol.\ = Solution-sourced.
  U = Uncertainty level; I = Impact level.
  Ratings: VL = Very Low, M = Medium, H = High, VH = Very High.}
\label{tab:hypotheses-overview}
\centering\small
\begin{tabular}{llllll}
\toprule
\textbf{ID} & \textbf{Short label} & \textbf{Team} & \textbf{Source} &
\textbf{U} & \textbf{I} \\
\midrule
H01 & Kotlin/JOOQ/Spring template readiness & Product  & Req. & VH & VH \\
H02 & Kafka direct exposure via platform framework & Platform & Req. & VL & M  \\
H03 & Decouple eng.\ logs from security/compliance audit & Platform & Req. & M  & VH \\
H04 & Blob storage GDPR right-to-erasure & Product  & Req. & M  & VL \\
H05 & Credit-bureau context-agnostic restructuring & Product  & Sol. & H  & VH \\
H06 & Billing aggregation chunked keyset pagination & Product  & Sol. & VL & VH \\
H07 & Blob ownership transfer across service domains & Product  & Req. & VL & VH \\
H08 & Blob cross-client sharing governance & Product  & Req. & VL & VH \\
H09 & Dead Letter Topic (DLT) for audit vs.\ offset-based replay & Platform & Sol. & M  & VH \\
H10 & Synchronous retry for transient consumer failures & Platform & Sol. & M  & H  \\
H11 & Failure state inside vs.\ outside Kafka flow & Platform & Sol. & M  & H  \\
H12 & Error classification before retry & Platform & Sol. & M  & H  \\
H13 & Payload immutability before re-enqueueing & Platform & Sol. & VL & VH \\
\bottomrule
\end{tabular}
\end{table}

\subsection{Source Types}
\label{subsec:source-types}

Eight of the thirteen hypotheses were classified as \textit{Requirements}-sourced:
the uncertainty arose from an incompletely specified quality-attribute goal or
an external obligation whose implementation path is unknown. Five were classified
as \textit{Solution}-sourced: the team has an existing or candidate architectural
decision and is uncertain about its consequences. The predominance of
requirements-sourced hypotheses on the Product~Team (five out of six) and
the predominance of solution-sourced hypotheses on the Platform~Team's Kafka
cluster (four out of five Kafka hypotheses) reflects the different nature of
each team's work: the Product~Team faces open questions about whether architectural
goals can be met at all, while the Platform~Team is assessing the trade-offs of
specific implementation choices for a shared infrastructure component.

\subsection{Quality Attribute Distribution}
\label{subsec:qa-distribution}

Table~\ref{tab:qa-distribution} summarizes the frequency of quality attributes
across the thirteen hypotheses. Attributes were assigned by the hypothesis authors
using the multi-select form in Hypo-Stage; a hypothesis may carry more than one
attribute. \emph{Because a single hypothesis was frequently classified under more
than one quality attribute (NFR), the counts in Table~\ref{tab:qa-distribution} sum
to well more than thirteen}: the fourteen attributes together account for thirty-seven
classifications across the thirteen hypotheses.

\begin{table}[htbp]
\caption{Quality attribute frequency across the thirteen hypotheses.}
\label{tab:qa-distribution}
\centering\small
\begin{tabular}{lcc}
\toprule
\textbf{Quality Attribute} & \textbf{Count} & \textbf{\% of hypotheses} \\
\midrule
Reliability        & 6 & 46\% \\
Auditability       & 5 & 38\% \\
Performance        & 4 & 31\% \\
Interoperability   & 4 & 31\% \\
Extensibility      & 3 & 23\% \\
Usability          & 3 & 23\% \\
Scalability        & 2 & 15\% \\
Modifiability      & 2 & 15\% \\
Flexibility        & 2 & 15\% \\
Configurability    & 2 & 15\% \\
Maintainability    & 1 &  8\% \\
Security           & 1 &  8\% \\
Observability      & 1 &  8\% \\
Reusability        & 1 &  8\% \\
\bottomrule
\end{tabular}
\end{table}

Reliability and Auditability emerge as the two dominant concerns. The prominence
of Auditability is especially notable and directly traceable to two contextual
forces: the organization's internal security audit requirements, which drove H03,
and the applicable privacy regulation governing document storage (H04),
as well as the Kafka immutability and DLT governance questions (H09, H13).
This concentration reveals a recurring archetype in which external compliance
obligations materialized into architectural hypotheses — a pattern discussed in
Section~\ref{sec:group-b-patterns} as the \textit{Regulatory Trigger Hypothesis}.

\subsection{Uncertainty and Impact Profile}
\label{subsec:uncertainty-impact}

Nine of the thirteen hypotheses carry a \textit{Very High} or \textit{High}
impact rating, signalling that the teams were not registering peripheral concerns
but uncertainties with potentially significant consequences for schedule, rework,
or system quality. Uncertainty ratings are more evenly distributed: five
hypotheses are rated \textit{Very Low} (the team has strong evidence but the
hypothesis is tracked for governance reasons), four are \textit{Medium}, one
is \textit{High}, and one \textit{Very High}.

The combination of Very Low uncertainty and Very High impact (H06, H07, H08,
H13) is particularly instructive. In ArchHypo's model, this quadrant maps to
hypotheses that \textit{can be postponed} — the team already has sufficient
evidence to act, so the immediate risk is low — yet the downstream consequences
of acting incorrectly are severe. Hypo-Stage's prioritization dashboard surfaced
these hypotheses in the "Can Postpone" segment of its focus panel, enabling
teams to consciously defer them without losing visibility. This represents
a practical instance of the \textit{Postpone Uncertain Decisions} pattern
described in the ArchHypo pattern literature \citep{silva2024patterns}.

\subsection{Technical Plans in Use}
\label{subsec:technical-plans}

Seven of the thirteen hypotheses had at least one associated technical plan at
the end of the observation period. Together these seven hypotheses carried thirteen
individual technical planning entries. The complete catalog of plan titles and
strategies is reproduced in Section~\ref{sec:app-technical-plans} of Appendix~B.
Table~\ref{tab:plan-strategies} summarizes
the strategies used.

\begin{table}[htbp]
\caption{Technical plan strategies observed across the hypothesis data.}
\label{tab:plan-strategies}
\centering\small
\begin{tabular}{lcc}
\toprule
\textbf{Strategy} & \textbf{Plans} & \textbf{Hypotheses using it} \\
\midrule
Guideline Implementation & 3 & 3 \\
Spike (Architectural Spike) & 3 & 2 \\
Modularization / Flexibility & 2 & 1 \\
Experiment & 2 & 2 \\
Tracer Bullet & 1 & 1 \\
Software Analytics (Trigger) & 1 & 1 \\
Other (implementation) & 1 & 1 \\
\midrule
\textbf{Total} & \textbf{13} & — \\
\bottomrule
\end{tabular}
\end{table}

Guideline Implementation and Architectural Spikes were the most frequently
chosen technical plans, consistent with findings from the PropSEP/Jetsoft evaluation
\citep{silva2025archhypo}. The Tracer Bullet plan appeared once for hypothesis H10,
where the Platform~Team planned a controlled injection of transient failures into
a real Kafka consumer path to measure ordering and latency characteristics. This
application illustrates how the Tracer Bullet differs from an Architectural Spike. Rather than
an isolated experiment, the Tracer Bullet plan exercises the full production consumer path
under representative conditions \citep{silva2024agilepatterns}.

Six hypotheses had no technical plan created within the observation window. Three
of these (H01, H04, H02) were registered but not yet acted upon; two (H07, H08)
were validated by planning actions taken \textit{before} the observation window
began, and their technical plans were entered into Hypo-Stage only afterward, as
a record of decisions already made rather than plans that guided work during the
study; and one (H05) was under active discovery, with its plan document still
being prepared.

\subsection{Lifecycle and Status}
\label{subsec:lifecycle-status}

At the end of the observation period, ten hypotheses had \textit{Open} status,
one was \textit{In Review} (H05), and two were \textit{Validated} (H07 and H08,
both in the document storage domain). No hypothesis was invalidated or discarded
during the study window. The predominance of the Open status is consistent with
the short duration of the study over the observation window
(Section~\ref{subsec:registration-dynamics}) and with the
nature of the uncertainties: most concern mid- to long-horizon architectural
concerns that teams do not expect to resolve within a single sprint cycle.

% ─────────────────────────────────────────────────────────────
\section{RQ3 — Perceived Benefits and Challenges}
\label{sec:rq3-results}
\label{sec:results-rq3}
% ─────────────────────────────────────────────────────────────

\subsection{Perceived Benefits}
\label{subsec:benefits}

Active listening sessions with all four interviewed participants (the two
technical leaders from each team) converged on four recurring themes regarding
the perceived value of using Hypo-Stage.
Because the sessions were not recorded
(Section~\ref{subsec:active-listening}), the attributions below are the
researcher's paraphrased reconstructions from dated field notes and
post-session written recollection---not verbatim quotations. The coding basis
for the perceived-benefit theme appears in Table~\ref{tab:code-traceability}
(\textit{tool-positive-catalog-integration}); challenge themes are developed in
Section~\ref{subsec:challenges}.

\paragraph{Externalization of tacit architectural knowledge:}
Both teams described the act of registering a hypothesis as valuable in itself.
Participants noted that seeing a hypothesis written down with a clear statement,
an explicit uncertainty rating, an impact rating, and a related quality attribute
made it easier to share the concern with other colleagues who had not been involved
in the original discussion.

\paragraph{Prioritization support:}
The \textit{Where to Focus} panel of the Hypo-Stage dashboard, distinguishing
hypotheses that \textit{Need Attention} (high uncertainty and high impact) from
those that \textit{Can Postpone} (low impact), was described by both technical
leaders as a useful input in architectural conversations. One technical leader
recalled that the ``Can Postpone'' category in particular gave them a way to
explain to other engineers outside of the team why a known concern was being
deliberately deferred rather than being ignored.

\paragraph{Collective memory for architectural decisions:}
Both teams observed that the hypothesis backlog served a documentation function
that was distinct from Architecture Decision Records (ADRs) or design documents.
Whereas ADRs capture decisions after they are made, an architectural hypothesis
captures an open question before a decision is available. Participants described
this as filling a gap in their current toolchain: they had good mechanisms for
recording what was decided but poor mechanisms for recording what was still
unknown.
Chapter~\ref{chap:conclusion} revisits this distinction when discussing implications for practice.

\paragraph{Catalog and Backstage workflow integration:}
Both teams described surfacing Hypo-Stage within the Backstage Internal Developer
Portal as a perceived benefit rather than merely an adoption convenience.
Because Backstage already served as the daily entry point for service catalog,
documentation, and CI/CD status, engineers could review and register hypotheses
without navigating to a separate tool or breaking out of the component context
they were already working in. Both technical leaders recalled that the
EntityHypothesesTab, which lists only hypotheses linked to the catalog component
currently being viewed, reduced the cognitive cost of context-switching into
hypothesis management. The open code \textit{tool-positive-catalog-integration}
labels this reconstructed theme; Section~\ref{subsec:refinement-backstage}
later describes how the study operationalized it as a refinement.

\subsection{Perceived Challenges}
\label{subsec:challenges}

Three challenges were consistently raised across participants. As with the
benefits above, these are reconstructed, role-attributed observations rather
than verbatim quotations; their open-code anchors are
\textit{falsifiability-confusion} and
\textit{registration-barrier-time-pressure}
(Table~\ref{tab:code-traceability}).

\paragraph{Negative-formulation learning curve:}
In this study, the teams and the researcher adopted a local convention of
registering only hypotheses phrased to specify the evidence under which the team
would conclude the hypothesis is incorrect---an inversion of the team's belief
into a ``what would prove this false'' form. All four participants described this
inversion as non-intuitive: engineers are accustomed to stating what they believe
to be true, and rephrasing that belief into its refutation condition requires a
deliberate cognitive effort that is easy to skip under time pressure. This
convention is stricter than what ArchHypo actually prescribes: the technique asks
only that a hypothesis be a \emph{falsifiable statement} (a refutable, testable
claim in which the possibility of proving it false exists, formulated as a simple,
direct statement of the team's current belief; Section~\ref{subsec:archhypo-core})
\citep{silva2024tse}. The cognitive cost observed here is therefore largely a
property of the study's negative-formulation convention rather than of ArchHypo's
guidance. The underlying \emph{learning curve}, however, is consistent with the
difficulty reported in both prior ArchHypo evaluations \citep{silva2024tse,
silva2024ieeesw}; Section~\ref{sec:trustworthiness} discusses this deviation as a
threat to construct validity. The open code
\textit{falsifiability-confusion} labels this reconstructed theme.

\paragraph{Registration time cost:}
Both technical leaders described the difficulty of dedicating the focused time
required to register a hypothesis properly. Fast-moving architectural micro-decisions, 
usually resolved within a day or less, were rarely registered because, by the time a
hypothesis could be written, the question had already been answered informally.
This creates a systematic under-representation of short-lifecycle uncertainties
in the hypothesis backlog, which may be consequential when the same question
re-emerges in a future context. The open code
\textit{registration-barrier-time-pressure} labels this reconstructed theme.

\paragraph{Absence of team-visible ceremonies:}
Neither team had established a recurring ceremony, such as an architectural
review or a hypothesis triage meeting, in which Hypo-Stage records were reviewed
as a team. Tool usage remained largely individual during the observation period.

% ─────────────────────────────────────────────────────────────
\section{RQ4 — Refinements to Hypo-Stage}
\label{sec:rq4-results}
\label{sec:results-rq4}
\label{sec:refinements}
% ─────────────────────────────────────────────────────────────

This section presents the refinements implemented in Hypo-Stage from commit
\texttt{ceee509} (the artifact boundary established by \citealt{correa2025}) to
commit \texttt{3a53d4e}. Each refinement was traced
to a concrete observation from the study and linked to the research question it
primarily addresses. All refinements address RQ4 by definition; selected
refinements also related to RQ1 (adoption ease) and RQ3 (usability and learning
curve).

\subsection{Backstage Integration and Catalog Surfacing}
\label{subsec:refinement-backstage}

\textbf{Observation:} During the training session and the first week of
observation, both teams navigated to Hypo-Stage via a standalone entry point
that was separate from their regular Backstage workflow. Engineers who were not
technical leaders rarely revisited the tool after the initial session.

\textbf{Refinement:} Hypo-Stage was wired as a first-class Backstage experience,
with routable pages for home, hypothesis creation, detail view, and edit flows.
Crucially, a catalog component tab (\texttt{EntityHypothesesTab}), shown in Figure~\ref{fig:entity-hypotheses-tab}, was introduced,
surfacing all hypotheses linked to the component an engineer is currently viewing
in the Backstage catalog.

\textbf{Impact:} Reducing navigation friction lowered the perceived cost of
checking the hypothesis backlog while already working with a service's catalog
entry.

\subsection{Backend API and Filtering}
\label{subsec:refinement-api}

\textbf{Observation:} Engineers with cross-team or platform-wide responsibilities
needed to filter hypotheses by team and by specific service. The initial
version of Hypo-Stage returned all hypotheses without server-side filtering,
which degraded usability as the hypothesis data grew.

\textbf{Refinements:} Three new API capabilities were implemented:
\begin{itemize}
  \item \texttt{GET /hypotheses} extended with \texttt{entityRef} and \texttt{team}
    query parameters;
  \item \texttt{GET /hypotheses/referenced-entity-refs} to populate the component
    filter without listing the entire catalog;
  \item \texttt{GET /hypotheses/teams} deriving team names from catalog metadata
    for components referenced by hypotheses.
\end{itemize}

\textbf{Impact:} On the Hypo-Stage home page, these capabilities enabled filtering
of the hypothesis list by team and by service (the Backstage catalog entity linked
to a hypothesis), so engineers with cross-team or platform-wide responsibilities
could narrow the backlog without loading the full unfiltered dataset.

\subsection{Prioritization Dashboard}
\label{subsec:refinement-dashboard}

\textbf{Observation:} The technical leaders expressed the desire to quickly
identify which hypotheses deserved immediate attention and which could safely be
deferred. Manually scanning the full hypothesis list imposed an unnecessary
cognitive burden, especially when the hypothesis data grew.

\textbf{Refinement:} A \texttt{HypothesisListDashboard} component was introduced
above the hypothesis table (Figure~\ref{fig:prioritization-dashboard}), presenting:
\begin{itemize}
  \item Overview cards with total count, validated count, and hypotheses created
    in the last thirty days;
  \item A \textit{Where to Focus} section distinguishing \textit{Need Attention}
    (high uncertainty and high impact) from \textit{Can Postpone} (low impact),
    with explanatory subtext aligned with ArchHypo's prioritization model
    \citep{silva2025archhypo};
  \item Distribution panels for uncertainty and impact Likert values.
\end{itemize}
The dashboard respects entity, team, and component scoping.

\textbf{Impact (RQ3).} The "Can Postpone" segment was specifically cited by
participants as resolving a communication challenge: it provided a principled
vocabulary for explaining deliberate deferral to non-technical stakeholders.

\begin{figure}[ht]
  \centering
  \includegraphics[width=0.85\textwidth]{prioritization-dashboard}
  \caption{Prioritization Dashboard in Hypo-Stage: the focus panel surfaces hypotheses
    according to the urgency signals derived from their uncertainty--impact combination,
    operationalizing the Postpone Uncertain Decisions pattern from the ArchHypo pattern literature \citep{silva2024patterns}.}
  \label{fig:prioritization-dashboard}
\end{figure}

\subsection{Hypothesis Lifecycle Consistency}
\label{subsec:refinement-lifecycle}

\textbf{Observation:} In the initial implementation, create and update operations
were not wrapped in database transactions, which created a risk of partial writes
(hypothesis row updated without the corresponding event entry, or vice versa).
Artifact analysis revealed two records with missing event entries, suggesting that
this race condition had occurred in practice.

\textbf{Refinements:} Create and update paths were refactored to use explicit
database transactions. Delete was extended to cascade to technical planning rows
and events, with clear API semantics and a confirmation dialog that requires the
exact hypothesis statement before proceeding (paste-enabled to avoid penalizing
users with long statements). Success and error toast notifications were added
for destructive operations.

\subsection{Technical-Planning-Driven Assessment Updates}
\label{subsec:refinement-tp-assessment}

\textbf{Observation:} Engineers reported that after completing a technical plan
action, such as an architectural spike, they updated the hypothesis's uncertainty
rating separately from the plan record. This two-step process was error-prone
and led to hypotheses whose assessment was stale relative to their technical
plans.

\textbf{Refinements:} Technical planning create and update operations were
extended to accept optional \texttt{uncertainty} and \texttt{impact} fields.
When present, the service updates the parent hypothesis's Likert values in the
same transaction. The event model was refined to distinguish
\texttt{TECHNICAL\_PLANNING\_CREATE} and \texttt{TECHNICAL\_PLANNING\_UPDATE}
events from generic hypothesis updates, preserving a traceable audit trail of
which assessment changes were driven by plan actions.

\subsection{Evolution Chart Improvements}
\label{subsec:refinement-chart}

\textbf{Observation:} The initial uncertainty/impact evolution chart used opaque
numeric y-axis ticks that participants could not interpret without consulting the
Likert scale documentation separately. Multiple events on the same calendar day
produced ambiguous x-axis labels.

\textbf{Refinements:} The chart was updated to render verbal y-axis ticks
(Very Low through Very High). Event source was made visually distinguishable:
circles for manual or creation-driven points, squares for technical-planning-driven
points. A legend explaining the marker shapes was added. When multiple events
occur on the same calendar day, the x-axis uses date and time to avoid ambiguity.
Tooltips were improved to label technical-planning-driven assessment moves with
the plan number when identifiable.

\begin{figure}[ht]
  \centering
  \includegraphics[width=0.8\textwidth]{evolution-chart}
  \caption{Evolution Chart in Hypo-Stage: tracking uncertainty and impact rating
    changes across successive assessments for a registered hypothesis. The screenshot
    is illustrative---a representative example of the chart rather than a depiction
    of one specific hypothesis from the study dataset. Each point corresponds to a
    saved assessment; the chart makes visible whether technical plan actions are
    reducing uncertainty or impact over time.}
  \label{fig:evolution-chart}
\end{figure}

\subsection{Form Usability Improvements}
\label{subsec:refinement-forms}

\textbf{Observation:} Participants found the hypothesis creation form cognitively
demanding, particularly the uncertainty and impact selectors, which displayed
only numeric values, and the quality attribute field, which accepted free text
without suggesting valid values.

\textbf{Refinements:} Uncertainty and impact selectors were redesigned as Likert
star pickers showing both a numeric button (1–5) and an adjacent text label
(Very Low through Very High) with disambiguating icons. Quality attributes were
replaced by a multi-select autocomplete showing attribute name and description
in the dropdown. Entity references for linking hypotheses to catalog components
were given a catalog-backed autocomplete. Documentation link fields were given an
improved list editor with add/remove and basic URL validation. 
Figure~\ref{fig:hypothesis-form} shows the improved hypothesis creation form.

\begin{figure}[ht]
  \centering
  \includegraphics[width=0.75\textwidth]{hypothesis-form}
  \caption{Hypothesis creation form in Hypo-Stage after the usability-oriented form
    refinements in this section: inline contextual guidance adjacent to the
    Statement field scaffolds the formulation of falsifiable hypotheses,
    reducing the cognitive barrier at the moment of registration.}
  \label{fig:hypothesis-form}
\end{figure}

\subsection{Summary: Refinements Mapped to Observations}
\label{subsec:refinement-summary}

Table~\ref{tab:refinements-map} summarizes the nine categories of refinements,
the primary observation that motivated each, and the research question each
addresses most directly.

\begin{table}[htbp]
\caption{Refinements to Hypo-Stage mapped to motivating observations and
  research questions. RQ columns indicate primary (P) or secondary (S) linkage.}
\label{tab:refinements-map}
\centering\small
\setlength{\tabcolsep}{4pt}
\begin{tabular}{p{3.8cm}p{4.5cm}ccc}
\toprule
\textbf{Refinement} & \textbf{Motivating observation} &
\textbf{RQ1} & \textbf{RQ4} & \textbf{RQ3} \\
\midrule
Backstage catalog tab
  & Infrequent revisit by non-leaders outside native workflow
  & S & P & S \\
Backend API filtering
  & Cross-team users needed scoped views as hypothesis data grew
  & S & P & — \\
Prioritization dashboard
  & Manual scanning of full backlog was cognitively costly
  & — & P & S \\
Lifecycle transaction consistency
  & Partial writes produced stale event records
  & — & P & — \\
TP-driven assessment updates
  & Assessment and plan records drifted apart post-spike
  & — & P & S \\
Evolution chart improvements
  & Numeric y-axis and same-day collisions caused confusion
  & — & P & S \\
Form usability improvements
  & Numeric-only selectors; free-text QA field caused errors
  & — & P & P \\
Safer deletion UX
  & Long statements made confirmation hostile without paste
  & — & P & — \\
Toast notifications
  & No feedback after destructive operations caused re-submission
  & — & P & — \\
\bottomrule
\end{tabular}
\end{table}

% ─────────────────────────────────────────────────────────────
\section{Emerging Findings from the Case Study}
\label{sec:emerging-patterns}
% ─────────────────────────────────────────────────────────────

Beyond the direct answers to RQ1–RQ4, thematic analysis of observation notes,
reconstructed active-listening notes, and the combined hypothesis data produced two
families of emerging findings. The \textbf{Adoption-Condition Candidates}
describe team-behavior conditions that shape whether hypothesis engineering
becomes a sustainable practice; they are candidate conditions rather than
validated patterns. The \textbf{Structural Hypothesis Patterns} describe
emerging patterns in the content of the architectural hypotheses
themselves. These two families address orthogonal dimensions: the
Adoption-Condition Candidates ask \textit{when} hypothesis engineering works;
the Structural Hypothesis Patterns ask \textit{what shape} hypotheses take in
practice. Table~\ref{tab:emerging-patterns} at the end of this section
summarizes the nine items across both families.

\begin{quote}
\textbf{Scope and statistical caveat:}
The findings described in this section were identified through thematic analysis
of thirteen architectural hypotheses produced by two software engineering teams
within a single organizational context over the observation window
(Section~\ref{subsec:registration-dynamics}).
This dataset is sufficient to motivate conceptual characterization and to propose
initial vocabulary for the registered hypothesis, but it does not support statistical
claims about frequency, generalizability, or prevalence.
The Adoption-Condition Candidates and Structural Hypothesis Patterns should
therefore be read as
\emph{emerging propositions} --- grounded in real empirical data and theoretically
coherent --- but requiring replication across additional teams, organizations, and
domains before they can be treated as consolidated findings.
Future work directions for each item are discussed in Section~\ref{sec:future-work}.
\end{quote}

% ─────────────────────────────────────────────────────────────
\subsection{Adoption-Condition Candidates for Hypothesis Engineering}
\label{sec:group-a-patterns}
% ─────────────────────────────────────────────────────────────

The four Adoption-Condition Candidates emerge by cross-referencing RQ1 and RQ3
findings. They describe organizational and individual-level conditions that the
study observed to facilitate or obstruct the adoption of hypothesis engineering
as a team practice. They extend the existing ArchHypo patterns observed by
\citet{silva2024patterns, silva2023initialpatterns} by addressing the social and
operational conditions in which ArchHypo mechanics operate in practice.

\emph{What these items are (and are not).} The four items described below
(A1--A4) are \emph{candidate conditions} under which hypothesis engineering can be sustained as a team practice. These are \emph{not} validated patterns. Unlike a pattern, which pairs a recurring
problem with a proven solution, each candidate here characterizes only the
\emph{problem and context} it was observed in. They are emergent rather than
confirmatory findings: each is supported by convergent evidence from both teams,
but they have not been tested in the broader population of software engineering
organizations. Accordingly, the \emph{solution component is deliberately
absent}: the derivation of actionable solutions from each candidate is deferred
to the future research agenda (Section~\ref{sec:future-work}), which
outlines the empirical validation required before any candidate could be
consolidated into a pattern or adopted as a guideline.

\subsubsection{Senior-Curated Hypothesis Backlog}
\label{subsubsec:pattern-a1}

\textbf{Problem:} In a team whose engineers operate at different scopes and
timescales, who is positioned to surface and sustain the registration of 
architectural hypotheses in mid- and long-terms?

\textbf{Observed evidence:} Twelve of the thirteen registered hypotheses were
authored by senior technical leaders — engineers who carried explicit
architectural responsibility and whose work extended beyond the current
sprint. Non-senior engineers did not sustain independent hypothesis registration
after the initial training episode. The hypotheses that were registered shared
a common trait: they concerned architectural uncertainties whose consequences would materialize
weeks or months in the future, making them invisible to typical sprint-scope
review practices.

\subsubsection{Long-Horizon Registration Priority}
\label{subsubsec:pattern-a2}

\textbf{Problem:} Many architectural micro-decisions are made and resolved within
hours or a single day. Rather than treating their speed as a blind spot to be
closed, the actionable observation is the inverse: fast-resolved micro-decisions
rarely repay the cost of registration, whereas the uncertainties genuinely worth
registering are the slow, high-consequence ones whose architectural effects only
materialize weeks or months later. Registration effort is therefore best
concentrated on long-horizon uncertainties rather than spread evenly across every
transient decision.

\textbf{Observed evidence:} Both technical leaders reported that the majority of
architectural decisions encountered day-to-day were resolved before a hypothesis
could be written, and that registration felt retrospectively unnecessary for them.
Notably, these fast-resolved decisions were not captured as Architecture Decision
Records (ADRs) either, so their absence from Hypo-Stage reflects a genuine
registration threshold rather than a tool-specific gap. Yet both also described
recurring situations where the same question re-surfaced in a new context,
requiring the team to reconstruct reasoning that had previously been resolved
informally and in a rapid pace --- a residual cost the leaders judged acceptable
relative to the effort of registering every micro-decision, which reinforces that
the payoff of registration concentrates in the slow, high-consequence uncertainties.

\subsubsection{Registration as a Slack-Dependent Practice}
\label{subsubsec:pattern-a3}

\textbf{Problem:} Translating a tacit architectural concern into a properly
formulated ArchHypo hypothesis requires focused, uninterrupted attention.
Sprint delivery pressure makes this kind of attention scarce and unreliable as
a basis for sustained practice.

\textbf{Observed evidence:} Senior leaders explicitly described the difficulty
of "dedicating proper time to register all possible hypotheses they encounter
day-to-day." The observed set of thirteen hypotheses is known to
under-represent the total uncertainties that leaders were aware of. Registration
activity clustered around protected and restricted discussion time rather than being distributed
across normal sprint work.

\subsubsection{Negative-Formulation Barrier}
\label{subsubsec:pattern-a4}
\textbf{Problem:} Under the negative-formulation convention adopted in this study
---registering hypotheses phrased to describe what would prove them false---writing
a hypothesis required a cognitive inversion of how engineers typically articulate
their beliefs. This inversion was a persistent adoption barrier even for
practitioners who understood the technique conceptually. The barrier is conditioned
on that convention: it is not an inherent requirement of ArchHypo, which prescribes
free-form falsifiable statements of the team's current belief \citep{silva2024tse}
and admits optional interrogative templates such as QUESt \citep{melegati2019quest}.

\textbf{Observed evidence:} All four interviewed participants mentioned that
formulating hypotheses under this convention was non-intuitive and difficult to
apply under time pressure. The associated \emph{learning curve} replicates findings
reported in the PropSEP/Jetsoft evaluation \citep{silva2025archhypo}; the
\emph{negative-formulation} framing, however, is specific to this study's convention.
Whether ArchHypo's prescribed free-form falsifiable statements avoid the barrier is
a hypothesis for future work rather than an established property of the technique.

% ─────────────────────────────────────────────────────────────
\subsection{Structural Hypothesis Patterns for Hypothesis Engineers}
\label{sec:group-b-patterns}
% ─────────────────────────────────────────────────────────────

The prior work in the ArchHypo research line has established the technique's
core concepts \citep{silva2025archhypo}, documented patterns for hypothesis
assessment, uncertainty management, and decision deferral
\citep{silva2024patterns, silva2023initialpatterns}, and reported on the
technique's learning curve and process-integration challenges
\citep{silva2024archhypo}. What this body of work has not yet addressed is
a question that only becomes answerable through empirical data on
hypotheses registered in practice:

\begin{quote}
\textit{What recurring shapes do architectural hypotheses take when teams apply
ArchHypo to real systems? And what structural properties of a hypothesis, its
source type, quality-attribute profile, and relation to its organizational
context, should an engineer recognize in order to write and manage it more
effectively?}
\end{quote}

This gap was also identified in calls for future work within the ArchHypo
pattern literature, which explicitly noted the need for "identifying and
documenting patterns for various types of hypotheses" and for "exploration of
common uncertainties related to specific quality attributes"
\citep{silva2024patterns}. The thirteen-hypothesis dataset produced by this
study constitutes the first empirical dataset that enables such identification.
The five patterns described below represent a \textbf{novel contribution}
of this work: they are grounded in real hypothesis data, not derived
from theoretical analysis, and they provide a vocabulary that 
engineers can use both to recognize familiar situations more quickly and to
communicate about hypothesis types across teams.

\emph{Statistical scope:} All five patterns were identified from the same
thirteen-hypothesis dataset. While each pattern maps to two or more
hypotheses in the dataset, the small dataset size means that no frequency claims are made. 
The contribution lies in naming and structuring recurring shapes so that future practitioners and researchers
have a vocabulary for recognition, not in establishing that any pattern
will appear with a particular frequency in other organizational contexts, nor
in prescribing solutions for each one. Solution-shaped guidance and
consolidation into the broader ArchHypo pattern language are delegated to the
future-work research agenda in (Section~\ref{sec:future-work}).

\subsubsection{Hypothesis Cluster}
\label{subsubsec:pattern-b1}

\textbf{Context:} A team is responsible for a shared platform component on which
multiple other service teams depend. An architectural question about that
component cannot be decomposed into a single hypothesis because the design space
has several interrelated dimensions, each of which involves a distinct trade-off.

\textbf{Problem:} Formulating a single hypothesis for a multi-dimensional
architectural concern either oversimplifies the problem (losing important
trade-off structure) or produces an unmanageably broad statement that is hard
to falsify.

\textbf{Observed evidence:} Five hypotheses (H09–H13) all concern the shared
Kafka streaming infrastructure and were registered on the same day, as the
Platform~Team's technical leader used Hypo-Stage to structure a design discussion.
The five hypotheses address: (i)~DLT audit semantics versus offset-based replay;
(ii)~synchronous versus asynchronous retry for transient consumer failures;
(iii)~failure-state placement inside versus outside the Kafka flow;
(iv)~error classification before retry; and (v)~payload immutability before
re-enqueueing. Each is independently falsifiable yet together they map the full
decision space for consumer-side failure handling in the shared streaming
platform.

\subsubsection{Regulatory Trigger Hypothesis}
\label{subsubsec:pattern-b2}
\textbf{Context:} A software system operates in a regulated domain such as financial
services, healthcare, or data privacy, where external compliance obligations
create architectural constraints whose implementation path is not fully known.

\textbf{Problem:} A compliance requirement imposes an architectural obligation
before the team has determined how to satisfy it. It is unclear whether the
existing architecture already satisfies the obligation, what changes would be
required if it does not, and what the cost of non-compliance would be.

\textbf{Observed evidence:} Two hypotheses in the dataset arise directly from
regulatory obligations: H04 (the GDPR, the General Data Protection Regulation, right-to-erasure requirement for binary
documents in the blob storage service) and H03 (organization-wide security
audit requirements for database log correlation across the platform framework).
In both cases the uncertainty is not about whether to comply — that is a legal
requirement that must be met — but about whether the current architecture \textit{can} comply and at
what technical cost.

\subsubsection{Service Decomposition Hypothesis}
\label{subsubsec:pattern-b3}
\textbf{Context:} A service has grown to handle multiple responsibilities such as domain
logic, external integrations, and identity management, and the team is
considering decomposing it into smaller services aligned with distinct domain
boundaries --- the recurring software-architecture-evolution challenge of
splitting a coarse-grained service into microservices along business-capability
or subdomain lines \citep{richardson2018microservices}. What the decomposition
would entail --- the correct contract boundaries and the migration path ---
remains unknown.

\textbf{Problem:} Deciding \textit{whether} and \textit{how} to decompose the service
is itself uncertain: the team does not know where the correct service boundary
should fall, which existing flows can be migrated without business-logic
changes, and what the downstream impact on other teams would be. As
\citet{richardson2018microservices} argues, decomposition can proceed by
business capability or by subdomain, but committing to a boundary without this
information risks a disruptive refactor that delivers less reuse than expected.

\textbf{Observed evidence:} H05 (credit-bureau service decomposition) captures this
precisely. The service currently handles identity, bureau integration, and
confidential document flows using domain-specific entities tied to a single
business vertical. The hypothesis asks whether these integrations can be
decoupled into a context-agnostic middleware service accepting only each provider's minimum
required identifiers, reusable across domains without duplicated cost. At the
end of the observation period, H05 was the only hypothesis in \textit{In Review}
status, with a discovery technical plan mapping (i)~the dependency matrix of
current versus real dependencies per integration, (ii)~an agnostic interface
contract proposal, and (iii)~a migration plan identifying which flows can isolate
immediately.

\subsubsection{Monitored Quality-Attribute Hypothesis}
\label{subsubsec:pattern-b4}
\textbf{Context:} A hypothesis targets a quality attribute --- such as
performance or scalability under load --- that is affected by many interacting
factors (data volume, access patterns, resource limits, runtime behavior). The
team has proposed a specific algorithmic or implementation approach and has
already conducted a capacity or volume test at scale, producing evidence that
the approach satisfies the target. However, the hypothesis has not been formally
closed because the implementation is still in progress.

\textbf{Problem:} The hypothesis occupies the lowest uncertainty quadrant
(Very Low) while retaining the highest impact rating (Very High), reflecting a
situation where the team \textit{knows} the solution works but the architectural
commitment is still consequential. Because the quality attribute depends on many
factors that can shift as the system evolves, it must be \textit{continuously
monitored} even after the uncertainty is validated: the hypothesis must be
tracked until implementation is complete and the attribute confirmed in
production, yet its characterization as "Open" may suggest unresolved doubt that
no longer exists.

\textbf{Observed evidence:} H06 (billing aggregation chunked keyset pagination)
exemplifies this. The team validated the approach — chunked keyset pagination
with isolated repository-level transactions — through a capacity test at five
times production scale (45 million rows total, 1.5 million rows per month).
Profiling confirmed that JPA stream processing caused first-level cache growth,
and that chunking scopes transactions to approximately ten thousand rows so the
garbage collector can reclaim memory, producing a healthy sawtooth heap pattern.
The hypothesis remains Open because the implementation is ongoing: the
uncertainty has been effectively resolved, but the performance quality attribute
continues to be monitored while the solution is rolled out.

\subsubsection{Shared-Resource Hypothesis}
\label{subsubsec:pattern-b5}
\textbf{Context:} A platform service manages a shared resource such as binary objects,
credentials, and configuration data, on which multiple product domains depend. New
business requirements demand new ways of accessing that resource across domain
boundaries, raising interacting consistency, security, and usability concerns
whose resolution is not yet known --- of which governance (ownership,
cross-domain read/write, and transfer between owning services) is one prominent
facet.

\textbf{Problem:} The team must decide how to model cross-domain access to the
shared resource in a way that keeps it consistent, secure, and usable across
domains. One central facet is governance --- specifically, which service "owns"
the resource at any given point, who may read or write it, and under what
conditions ownership can be transferred --- but these decisions interact with
broader consistency, security, and usability concerns. Without explicit
decisions, different product teams will adopt incompatible access patterns,
creating consistency violations that are difficult to detect and remedy.

\textbf{Observed evidence:} H07 (blob ownership transfer across service domains)
and H08 (blob cross-client sharing governance) are two instances of this pattern
in the secure blob storage domain. H08 was driven by a concrete usage scenario
in which a customer-facing channel receives documents that are later routed to
support tooling; the support team needs read access without breaking the blob's ownership invariants. 
Both hypotheses were validated through a Guideline technical plan composed of an alignment
session held before the observation window and registered in Hypo-Stage retrospectively,
producing a governance decision enabling prioritization of the transfer and sharing
implementation work.

\begin{table}[htbp]
\caption{Emerging findings identified in the case study: the four
Adoption-Condition Candidates (A1--A4) and the
five Structural Hypothesis Patterns (B1--B5). Both families describe problem and context only;
solution-shaped guidance for each item is deferred to the future research
agenda (Section~\ref{sec:future-work}).}
\label{tab:emerging-patterns}
\centering\small
\setlength{\tabcolsep}{4pt}
\begin{tabular}{clp{3.6cm}p{5.2cm}p{2.6cm}}
\toprule
\textbf{ID} & \textbf{Family} & \textbf{Name} & \textbf{Signature} &
\textbf{Primary evidence} \\
\midrule
A1 & Adoption candidate & Senior-Curated Hypothesis Backlog
   & Mid-/long-horizon uncertainties stay visible because senior engineers
     author them
   & 12 of 13 hypotheses senior-authored \\
A2 & Adoption candidate & Long-Horizon Registration Priority
   & Registration value concentrates in slow, high-consequence uncertainties;
     fast-resolved micro-decisions rarely repay the cost (and are absent from
     ADRs too)
   & Both technical leaders' testimony \\
A3 & Adoption candidate & Registration as a Slack-Dependent Practice
   & Well-formulated hypotheses require protected time, not sprint slack
   & Registration clustering at training and Kafka session \\
A4 & Adoption candidate & Negative-Formulation Barrier
   & The study's negative-formulation convention (not ArchHypo itself) makes
     hypothesis writing cognitively costly
   & All four interviewed participants \\
\midrule
B1 & Structural pattern & Hypothesis Cluster
   & Multi-dimensional platform concern decomposes into linked falsifiable
     hypotheses
   & H09--H13 (Kafka failure handling) \\
B2 & Structural pattern & Regulatory Trigger Hypothesis
   & External compliance obligation surfaces as tracked architectural
     uncertainty
   & H03, H04 \\
B3 & Structural pattern & Service Decomposition Hypothesis
   & Uncertainty about whether and how to decompose a service into microservices
     along domain boundaries
   & H05 (credit-bureau service decomposition) \\
B4 & Structural pattern & Monitored Quality-Attribute Hypothesis
   & Very-low uncertainty, very-high impact: quality attribute empirically
     validated but continuously monitored while implementation is open
   & H06 (billing aggregation) \\
B5 & Structural pattern & Shared-Resource Hypothesis
   & Cross-domain consistency, security, and usability concerns --- including
     ownership governance --- for a shared platform resource
   & H07, H08 \\
\bottomrule
\end{tabular}
\end{table}

% ─────────────────────────────────────────────────────────────
\section{Revisiting the Research Questions}
\label{sec:rq-revisit}
% ─────────────────────────────────────────────────────────────

This section concludes the analytical observations of the chapter by stating explicitly how
the preceding empirical and design material responds to the four research questions
formulated in Chapter~\ref{chap:introduction}. The intent is \emph{consolidative} where
no additional empirical claims are introduced. The aim is to make the logical
organization of the chapter legible relative to the research questions.

\subsection{Correspondence between research questions and chapter sections}

Table~\ref{tab:rq-chapter-map} summarizes where each research question is principally
addressed in this chapter, and where integrated interpretation (the
Adoption-Condition Candidates and Structural Hypothesis Patterns) extends
the principal sections.

\begin{table}[htbp]
\centering
\small
\setlength{\tabcolsep}{5pt}
\caption{Correspondence between the research questions (Chapter~\ref{chap:introduction})
  and the organization of evidence in this chapter.}
\label{tab:rq-chapter-map}
\begin{tabular}{p{0.12\linewidth} p{0.40\linewidth} p{0.40\linewidth}}
\toprule
\textbf{RQ} & \textbf{Principal analysis} & \textbf{Integrated interpretation} \\
\midrule
RQ1 & Section~\ref{sec:rq1-results} (tool use and adoption dynamics). &
Section~\ref{sec:group-a-patterns} (Adoption-Condition Candidates; read together with RQ3). \\
\midrule
RQ2 & Section~\ref{sec:rq2-results} (typological account of the data). &
Section~\ref{sec:group-b-patterns} (Structural Hypothesis Patterns). \\
\midrule
RQ3 & Section~\ref{sec:rq3-results} (benefits and challenges). &
Section~\ref{sec:group-a-patterns} (Adoption-Condition Candidates; read together with RQ1). \\
\midrule
RQ4 & Section~\ref{sec:rq4-results} (nine refinement categories, observation-led). &
--- \\
\bottomrule
\end{tabular}
\end{table}

\subsection{Brief interpretive synthesis}

Through the preceding analysis presented in this chapter, the four research questions tell a single story about how far a
tool-mediated instantiation of ArchHypo can carry a team, and where the remaining barriers lie.

\textbf{RQ1 and RQ3 converge on an organizational, not a tool, problem:}
Both questions produce evidence that Hypo-Stage lowered surface friction, 
facilitating context-switching, filtering, traceability. Participants
recognized the intellectual value of hypothesis management. Still, the same two
questions also show that authorship stayed concentrated in senior engineers,
who registered hypotheses clustered around protected discussion time rather than
distributing across sprint work, and that hypothesis writing under the study's
negative-formulation convention remained cognitively costly. The
Adoption-Condition Candidates (A1--A4) name the
conditions under which those findings hold: adoption breadth is limited by the
organization structure including senior curation, protected cognitive time,
statement-level support, and ceremonial touchpoints. These are conditions that a tool alone
cannot supply.

\textbf{RQ2 and the Structural Hypothesis Patterns reveal a content gap in the
existing ArchHypo pattern language:} The prior pattern literature
\citep{silva2024patterns, silva2023initialpatterns} names mechanics of
assessment, deferral, and uncertainty management, but does not name the
\emph{types} of hypotheses engineers actually register. When applied in a
compliance-heavy, platform-dominant environment, the thirteen registered
hypotheses concentrate on regulatory triggers, shared-resource governance,
context-agnostic restructuring, and multi-dimensional platform concerns, which are categories the
existing patterns do not index. The Structural Hypothesis Patterns therefore
complement the mechanics-focused patterns rather than competing with them, and provide
the first empirically grounded vocabulary for the content of hypotheses
that emerge in practice.

\textbf{RQ4 closes the formative-evaluation loop:} The nine refinement
categories are not a separate implementation narrative. Each maps to a
specific observed friction surfaced by RQ1, RQ2, and RQ3 questions: the form-level and
ceremony-related refinements track A3 and A4; the component-scoped tab
tracks B1-shaped concerns about multi-dimensional platform uncertainty; the
prioritization dashboard and API filtering respond to the adoption
dynamics observed under RQ1. Their derivation from the same observational
data is itself evidence that the design-science loop of observation,
refinement, and re-observation operated on the Hypo-Stage tool during the study.

Taken together, the chapter supports a bounded claim: tool-mediated
hypothesis engineering succeeds at the integration and traceability layers
and produces a useful vocabulary for the content of architectural
hypotheses, but its adoption ceiling lies at the organizational layer. The
Adoption-Condition Candidates are the natural hand-off to
Chapter~\ref{chap:conclusion}, where they are read as open empirical problems rather than as actionable guidelines.
